% Paper section fragment, not a standalone document.
% Requires amsmath, amssymb, amsthm and a theorem environment in the host.
% Integrated from the version-bound lower/upper proofs, 2026-10-07.
% No claim of completed M1 realizability, novelty audit, D2-a, or TAC readiness.

\section{Sharp information loss in the stationary AR(1) experiment}
\label{sec:ar1-sharp}

We use raw observations
\begin{equation}\label{eq:ar1-raw-model}
 Y_i=s_i b^h(t_i)+\mathbf1_{\{h=1\}}a_{t_i}+\epsilon_{t_i},
 \qquad s_i\in\{-1,1\},
\end{equation}
where the noise is common, known and stationary Gaussian with
$\operatorname{Cov}(\epsilon_j,\epsilon_\ell)=\sigma^2\rho^{|j-\ell|}$,
$\sigma>0$ and $|\rho|<1$. A fixed healthy prefix at task $+$ consists of
$n_0\ge1$ slots $j=-n_0+1,\ldots,0$. The known post-onset horizon is $H$;
$a_j=0$ on the prefix and $a_j=\rho_0+\eta j$ for $1\le j\le H$,
where $\rho_0\ge0$ and $\eta>0$. The calendar is deterministic and selected
before observing noise. Holding supplies one reading per slot. A change of
task excludes at least $k\ge1$ consecutive slots. Longer moves, consecutive
moves and empty round trips are allowed. There is no imposed terminal task.

The two hypotheses choose their own full healthy paths, each with rate at
most $r/2$ and unrestricted level. Their difference $z=b^0-b^1$ has the exact
retained-node class
\begin{equation}\label{eq:ar1-trace}
 K_\pi=\{z:|z_{i+1}-z_i|\le r(t_{i+1}-t_i)\},\qquad r>0.
\end{equation}
Indeed, any node vector in this class has a rate-$r$ piecewise linear
extension, and $z/2,-z/2$ realize the two original paths.
For $R_\pi=\{t_1<\cdots<t_N\}$, let $C_\pi$ be the principal submatrix of the
single stationary full covariance and $S_\pi=\operatorname{diag}(s_i)$.
No noise state is reset during missing observations. Define
\begin{align}\label{eq:ar1-information}
 G_H(\pi)&=\tfrac12\inf_{z\in K_\pi}
 (a_R-S_\pi z)^\top C_\pi^{-1}(a_R-S_\pi z),\nonumber\\
 O_H&=\tfrac12a_{\rm full}^\top C_{\rm full}^{-1}a_{\rm full},\qquad
 D_H^*=O_H-\max_{\pi\in\Pi_H}G_H(\pi).
\end{align}
These are Gaussian half-Mahalanobis informations. Under task demodulation,
both the mean and covariance must transform; retaining $C_\pi$ while changing
only the target does not represent \eqref{eq:ar1-raw-model}.

\begin{theorem}[Fixed-parameter AR(1) sharp deficit]\label{thm:ar1-sharp}
Fix $\rho_0\ge0$, $\eta,r>0$, $k\ge1$, $\sigma>0$, $|\rho|<1$ and a fixed
$n_0\ge1$. Put
\begin{equation}\label{eq:ar1-nu-beta}
 \nu=\frac{1-\rho}{\sigma^2(1+\rho)},\qquad
 \beta_k=(k+1)\nu-\frac{1-\rho^{k+1}}{\sigma^2(1+\rho^{k+1})}>0.
\end{equation}
Then
\begin{equation}\label{eq:ar1-main}
 D_H^*=\frac{\sqrt2}{5}\sqrt{\nu\beta_k r}\,\eta^{3/2}H^{5/2}
 +o(H^{5/2}).
\end{equation}
The symmetric terminal-balanced calendar described in the proof, with
$\xi_*=2\beta_k\eta/(\nu r)$, attains the coefficient.
For every fixed $\xi>0$ its finite-horizon deficit obeys
\begin{equation}\label{eq:ar1-calendar-upper}
 D_H(\pi_H^\xi)\le\frac1{10}
 \left(\frac{2\beta_k\eta^2}{\sqrt\xi}+\nu r\eta\sqrt\xi\right)H^{5/2}
 +O(H^2).
\end{equation}
This is a first-deficit-coefficient statement. It does not match the
$H^2$-order term or allow channel parameters to vary with $H$ toward the
stability or nondegeneracy boundary.
\end{theorem}

\begin{proof}
We prove a calendar-uniform converse and a full-history achieving dual.

\paragraph{Exact marginal information and missing-gap identities.}
Write $Q_R(x)=x_R^\top C_R^{-1}x_R$ and $q_R=Q_R/2$.
Iteration of the AR(1) recursion yields independent $d$-step innovation
sums of variance $\sigma^2(1-\rho^{2d})$, so
\begin{equation}\label{eq:ar1-markov}
 Q_R(x)=\frac{x_{t_1}^2}{\sigma^2}
 +\sum_{i=2}^N\frac{(x_{t_i}-\rho^{t_i-t_{i-1}}x_{t_{i-1}})^2}
 {\sigma^2(1-\rho^{2(t_i-t_{i-1})})}.
\end{equation}
An internal gap containing $\ell$ missing nodes, with span $m=\ell+1$,
replaces $m$ one-step terms by one $m$-step term. Hence its full-minus-marginal
quadratic loss is local and positive semidefinite; losses over internal gaps
add exactly even when gaps share a retained endpoint. At constant contrast
$A$ its squared-Mahalanobis loss is $\beta_\ell A^2$, with
\begin{equation}\label{eq:ar1-all-beta}
 \beta_\ell=(\ell+1)\nu-w_{\ell+1}/\sigma^2,\qquad
 w_m=\frac{1-\rho^m}{1+\rho^m}.
\end{equation}
For a local affine vector $a_j=A+\eta(j-c)$ centered between the retained
endpoints, reflection invariance of both full and endpoint marginal
precision makes the constant--slope cross term zero. The resulting loss is
\begin{equation}\label{eq:ar1-affine-gap}
 \Delta q_{\rm gap}(a)=\tfrac12\beta_\ell A^2+\tfrac12\eta^2E_\ell,
 \qquad E_\ell\ge0.
\end{equation}
For completeness, if $M$ denotes its quadratic loss matrix and $J$ reflects
the local indices, then $JMJ=M$, $J\mathbf1=\mathbf1$ and $Jh=-h$ for
$h_j=j-c$. Thus $\mathbf1^\top Mh=0$ and $E_\ell=h^\top Mh\ge0$.
Positivity follows also by completing the square after marginalizing the
missing coordinates.

Let $A_m=\sum_{j=0}^{m-1}\rho^j$ and $B_m=\sum_{j=0}^{m-1}\rho^{2j}$.
Algebra gives $\beta_\ell/\nu=m-A_m^2/B_m>0$, since equality in
Cauchy--Schwarz requires $\rho=1$ for $m\ge2$.
Nested deletion in a fixed larger complete sequence shows that
$\beta_\ell$ is nondecreasing in $\ell$. Since $\beta_\ell/\ell\to\nu>0$,
\begin{equation}\label{eq:ar1-bstar}
 b_*:=\inf_{\ell\ge k}\beta_\ell/\ell>0.
\end{equation}
At each post-onset gap its left endpoint has nonnegative affine amplitude.
Every missing amplitude is therefore at most twice its center amplitude.
Writing $\mathcal D_\pi$ for all missing nodes gives
\begin{equation}\label{eq:ar1-dead-budget}
 T:=2\sum_{\rm gaps}\Delta q_{\rm gap}(a)
 \ge\frac{b_*}{4}\sum_{j\in\mathcal D_\pi}a_j^2.
\end{equation}

\paragraph{Uniform prefix and terminal reduction.}
The stationary full precision has absolute row sums at most
$C_\rho=(1+|\rho|)/[\sigma^2(1-|\rho|)]$. Thus its norm is at most
$C_\rho$, and covariance principal-submatrix interlacing gives
$\|C_R^{-1}\|\le C_\rho$ for every calendar.
Extend the target affinely to the fixed prefix and temporarily write
$\widetilde O,\widetilde G,\widetilde D$ for the modified experiment.
The target changes in only $n_0$ fixed coordinates, by white norm $O(1)$.
Distance to a set is 1-Lipschitz, and $0$ is feasible, so its distance is
at most the target white norm $O(H^{3/2})$.
The squared-distance inequality therefore gives, uniformly in the calendar,
$\widetilde D-D=O(H^{3/2})$.
A terminal sequence of moves without a later reading is dominated by
continuing to observe the preceding task. Data processing for every common
full path, followed by an infimum over that class, preserves this order.
Thus the optimal calendar may be assumed to observe slot $H$.
With the retained prefix, every post-onset missing interval is internal.
Prefix gaps never arise, and the first possible gap has left index $0$ and
amplitude $\rho_0\ge0$, as used in \eqref{eq:ar1-dead-budget}.

\paragraph{One full-history adversarial path.}
Fix $\varepsilon\in(0,1)$ and split the late interval
$[\lceil\varepsilon H\rceil,H]$ into complete macroblocks of length
$L=\lfloor H^{3/4}\rfloor$. On each clipped observed run of length $n$ use
$z_i=s_i r\min(i,n-1-i)$, $i=0,\ldots,n-1$.
All missing nodes, run endpoints, macroblock boundaries, prefix and unused
tail have $z=0$. Piecewise-linear interpolation forms a single legal
rate-$r$ full path, not separate paths with reset levels.
Set $h=s z$ on reads and zero on missing slots.
It is nonnegative and globally rate-$r$, and $x=a-h$ equals $a$ on every gap
including its two retained endpoints.

Expanding \eqref{eq:ar1-markov} on the complete contiguous record gives
\begin{equation}\label{eq:ar1-full-energy}
 q_{\rm full}(x)=\tfrac\nu2\sum x_j^2
 +\frac{c_\rho}{2}(x_{\rm first}^2+x_{\rm last}^2)
 +\frac{d_\rho}{2}\sum(\Delta x_j)^2,
\end{equation}
where $c_\rho=\rho/[\sigma^2(1+\rho)]$ and
$d_\rho=\rho/[\sigma^2(1-\rho^2)]$.
The global endpoints of $h$ vanish, so the endpoint terms cancel between
$a$ and $a-h$. Their slope-energy difference is $O(H)$, uniformly over
the calendar, even when $\rho<0$.
Evaluating the above legal path and subtracting its exact gap losses yields
\begin{equation}\label{eq:ar1-adversarial}
 2\widetilde D_H(\pi)\ge
 \nu\sum_{R_\pi}(2a_jh_j-h_j^2)+T-O(H).
\end{equation}

\paragraph{Calendar-uniform macroblock balance.}
In a macroblock let $A$ be the minimum amplitude, $d$ the number of missing
slots, $m$ the number of clipped observed runs and $n_i$ their lengths.
Then $\sum n_i=L-d$. For large $H$, uniformly in late macroblocks,
$A=\Theta(H)$ and
$C=r(2A-rL/2)/4>0$.
The tent height is at most $rL/2$ and its area satisfies
$\sum_{j=0}^{n-1}\min(j,n-1-j)\ge(n^2-2n)/4$. Consequently
\begin{equation}\label{eq:ar1-tent-area}
 \sum_{\rm macro}(2ah-h^2)\ge C\sum_i n_i^2-2CL.
\end{equation}
For $m>0$, Cauchy--Schwarz gives $\sum n_i^2\ge(L-d)^2/m$.
The $m-1$ gaps between clipped runs lie wholly in that macroblock, have
length at least $k$ and center amplitude at least $A$. By
\eqref{eq:ar1-affine-gap} and monotonicity,
$T\ge\sum_{\rm macros}\beta_kA^2(m-1)$.
For $m=0$ this last term is negative and remains a valid weaker bound.
Equation~\eqref{eq:ar1-dead-budget} also gives
$T\ge(b_*/4)\sum_{\rm macros}A^2d$.
Gaps crossing macro boundaries are absent from the first bound but their
missing nodes are assigned once in the second.

Fix $\theta\in(0,1)$ and split the same global $T$ into
$(1-\theta)T+\theta T$ before applying these two bounds.
Put $W=(1-\theta)\beta_k$. For $m>0$, the macroblock contribution to
\eqref{eq:ar1-adversarial} is at least
\[
 \nu C(L-d)^2/m+WA^2m+\theta b_*A^2d/4-WA^2-2\nu CL.
\]
AM--GM bounds the first two terms below by
$2A\sqrt{\nu CW}(L-d)$. Because $C\le rA/2$,
$A\ge32\nu Wr/(\theta^2b_*^2)$ suffices for
$\theta b_*A^2/4\ge2A\sqrt{\nu CW}$.
This holds uniformly for all late macroblocks as $H\to\infty$.
It absorbs the missing-slot subtraction. If $m=0$ and $d=L$, the same
absorption applied directly to the dead term gives the identical bound
\begin{equation}\label{eq:ar1-macro-bound}
 2A\sqrt{\nu CW}\,L-WA^2-2\nu CL.
\end{equation}
This argument needs no run-length regularity, sparse-switch assumption or
upper bound on missing-gap length.

Now $C=(rA/2)(1+O(H^{-1/4}))$. The accumulated $WA^2$ terms are
$O(H^3/L)=O(H^{9/4})$, the $CL$ terms are $O(H^2)$, and the coefficient and
Riemann-sum errors are $O(H^{9/4})$. The final incomplete macroblock also
affects that Riemann sum by at most $O(H^{3/2}L)=O(H^{9/4})$.
All constants are calendar-independent for fixed $\varepsilon,\theta$.
Moreover
\[
 \sum_{\rm macros}A^{3/2}L=
 \frac25\eta^{3/2}(1-\varepsilon^{5/2})H^{5/2}+o(H^{5/2}).
\]
Taking the infimum deficit before the limit, returning from
$\widetilde D$ to $D$, and then sending $\varepsilon,\theta$ to zero proves
\begin{equation}\label{eq:ar1-converse-final}
 \liminf_{H\to\infty}D_H^*/H^{5/2}
 \ge\frac{\sqrt2}{5}\sqrt{\nu\beta_k r}\,\eta^{3/2}.
\end{equation}

\paragraph{Terminal-balanced achieving calendar.}
For fixed $\xi>0$, a symmetric block with $n=2m$ comprises $m$ reads at
$+$, $k$ missing slots, $n$ reads at $-$, $k$ missing slots and $m$ reads at
$+$. Its duration is $L=2(n+k)$ and center is $c=(L+1)/2$.
With budget $H-1$, fit the maximal number $M$ of blocks with base
$\bar n_\ell=2\lceil\xi\ell/2\rceil$. Distribute remaining multiples of four
evenly over these blocks, with each unit increasing $n_\ell$ by two and a
fixed left-to-right tie rule. Leave one to four final observed slots at $+$.
If no block fits, stay. Since base cumulative duration is $\xi M^2+O(M)$
and the next block has duration $O(M)$, each added $n_\ell$ is $O(1)$.
Thus, writing $x_\ell$ for its start,
\begin{equation}\label{eq:ar1-block-scales}
 M=\sqrt{H/\xi}+O(1),\quad n_\ell=\xi\ell+O(1),\quad
 a_\ell=\rho_0+\eta(x_\ell+c_\ell)=\eta\xi\ell^2+O(\ell+1).
\end{equation}
These are realizable physical blocks, concatenating at task $+$ without
additional movements. Reserving the final reading avoids a partially
completed long terminal block.

\paragraph{A full raw dual with exact signed moment.}
For any $\lambda$ with zero sum, the unrestricted-level rate class has support
\begin{equation}\label{eq:ar1-support}
 h_{K_\pi}(\lambda)=r\sum_i(t_{i+1}-t_i)
 \left|\sum_{j\le i}\lambda_j\right|;
\end{equation}
this follows by summation by parts and independent bounded increments.
For nonzero sum its support is infinite.
Let $P=C_R^{-1}$ for the single complete retained calendar, and construct a
global auxiliary raw vector $u$. On the prefix and final one to four reads
put $u=0$. On an active block, defined by
$a_\ell\ge r(k+n_\ell-1)/2$, put $u_i=a_\ell-s_i z_i^{\rm loc}$, with
$c_0=(k+1)/2$ and local levels
\begin{align}\label{eq:ar1-local-levels}
 z_i^{+,L}&=r(c_0+m-i),&1\le i\le m,\nonumber\\
 z_i^-&=-r[c_0+\min(i-1,n-i)],&1\le i\le n,\\
 z_i^{+,R}&=r(c_0+i-1),&1\le i\le m.\nonumber
\end{align}
On inactive blocks set $u=0$. Only a uniformly finite initial set can be
inactive because $a_\ell$ grows quadratically and its threshold linearly.
These local levels are a device for constructing a dual; they need not join
to a feasible global primal healthy path.

Each active block has $0\le u_i\le a_\ell$, equal gap endpoint values
$W_\ell=a_\ell-rc_0$, and $\sum s_i u_i=0$.
The local direction $s_i u_i$ has support
\begin{equation}\label{eq:ar1-local-support}
 h_\ell^0=\frac{a_\ell r}{2}n_\ell(n_\ell+2k)
 -4r^2\sum_{j=0}^{n_\ell/2-1}(c_0+j)^2.
\end{equation}
To check this, its partial sums are nonnegative before the middle of the
negative run and nonpositive afterward, returning to zero at that midpoint
and at the block end. Whenever a partial sum $Q_i\ne0$, the local nodes obey
$\Delta z_i^{\rm loc}=-r\Delta t_i\operatorname{sign}(Q_i)$.
The central flat edge has $Q_i=0$, and a gap edge has $\Delta t_i=k+1$.
Hence \eqref{eq:ar1-support} is attained. The fourfold multiplicity of each
level gives $\sum s_i z_i^{\rm loc}=rn(n+2k)/2$ and
$\sum(z_i^{\rm loc})^2=4r^2\sum(c_0+j)^2$, proving
\eqref{eq:ar1-local-support}.

Let $\mathbf e_i$ denote coordinate basis vectors. Expand
\eqref{eq:ar1-markov} to write the full retained precision as $\nu I$,
consecutive-read edge matrices
$d_\rho(\mathbf e_i-\mathbf e_j)(\mathbf e_i-\mathbf e_j)^\top$, a matrix
\[
 \begin{pmatrix}\mathcal A_m&-\gamma_m\\-\gamma_m&\mathcal A_m\end{pmatrix},\quad
 \mathcal A_m=\frac1{\sigma^2(1-\rho^{2m})}-\frac1{\sigma^2(1+\rho)},\quad
 \gamma_m=\frac{\rho^m}{\sigma^2(1-\rho^{2m})},
\]
at each actual gap with $m=k+1$, and only two global endpoint corrections
$c_\rho$. There are no stationary initializations at artificial block joins.
Put $\delta_m=\mathcal A_m-\gamma_m
=1/[\sigma^2(1+\rho^m)]-1/[\sigma^2(1+\rho)]$.
Every consecutive-read edge has equal task signs; every actual gap has
opposite signs and equal $u$ endpoint values. Since global first and last
$u$ are zero, the exact vector identity is
\begin{equation}\label{eq:ar1-exact-zero-moment}
 SPu=\nu Su+
 \sum_{\Delta t=1}s_i d_\rho(u_i-u_j)(\mathbf e_i-\mathbf e_j)
 +\sum_{\rm gaps}\delta_m W_\ell s_{\rm left}
 (\mathbf e_{\rm left}-\mathbf e_{\rm right}).
\end{equation}
Every correction is a zero-mass pair, so $\sum SPu=0$ exactly.
All statements hold for negative $\rho$ as well; estimates below take
absolute values where needed.

\paragraph{True cross-block covariance and support remainder.}
Choose the raw direction $v=Pu$. The quadratic Fenchel inequality implies
$G\ge v^\top a_R-v^\top Cv/2-h_{K_\pi}(Sv)$.
Here $v^\top Cv=u^\top Pu$ exactly, using the full retained covariance.
Writing $e=a_R-u$ and $O_{\rm obs}=a_R^\top Pa_R/2$ gives
\begin{equation}\label{eq:ar1-dual-deficit}
 D_H(\pi)\le O_H-O_{\rm obs}+e^\top Pe/2+h_{K_\pi}(SPu).
\end{equation}
The actual direction $v$ may have nonzero prefix and tail weights; zero
auxiliary values do not discard observations.
Support sublinearity, restriction of the full healthy class to each block,
and support $r\Delta t|c|$ for a pair $c(\mathbf e_i-\mathbf e_j)$ give
\begin{equation}\label{eq:ar1-support-bound}
 h_{K_\pi}(SPu)\le\nu\sum_{\rm active}h_\ell^0
 +r|d_\rho|\sum_{\Delta t=1}|u_i-u_j|
 +r(k+1)|\delta_m|\sum_{\rm gaps}|W_\ell|.
\end{equation}
Within blocks the total variation is at most $2r\sum n_\ell$, and all joins
add at most $2\sum a_\ell$. Also $\sum_{\rm gaps}|W_\ell|\le2\sum a_\ell$.
Thus the two correction terms are $O(H^{3/2})$.
Equation~\eqref{eq:ar1-local-support}, after discarding its negative square
term, gives
\begin{equation}\label{eq:ar1-support-leading}
 h_{K_\pi}(SPu)\le\frac{\nu r}{2}\sum a_\ell n_\ell^2+O(H^2),
\end{equation}
since $\sum a_\ell n_\ell=O(H^2)$. No equality of global and independently
optimized local supports has been assumed.

\paragraph{Affine remainder and exact oracle cost.}
On active blocks $e_i=\eta(t_i-c_\ell)+s_i z_i^{\rm loc}$, so
$|e_i|\le(\eta+r)(n_\ell+k)$. The precision norm bound established above
and the at most four tail errors give
\[
 e^\top Pe\le C_\rho\{2(\eta+r)^2\sum n_\ell(n_\ell+k)^2
 +4(\rho_0+\eta H)^2+O(1)\}=O(H^2).
\]
Inactive blocks contribute only $O(1)$, and prefix errors are zero.
The two gaps in a block have centers at distances $\pm(n_\ell+k)/2$ from
its center. By \eqref{eq:ar1-affine-gap}, their full oracle loss is exactly
\[
 O_H-O_{\rm obs}=\sum_\ell\left\{\beta_k a_\ell^2
 +\frac{\beta_k\eta^2}{4}(n_\ell+k)^2+\eta^2E_k\right\}
 =\beta_k\sum a_\ell^2+O(H^{3/2}).
\]
These are all actual internal post-onset gaps; there is no initial or terminal
missing record in the construction.
Substitution in \eqref{eq:ar1-dual-deficit} yields
$D_H(\pi_H^\xi)\le\beta_k\sum a_\ell^2+(\nu r/2)\sum a_\ell n_\ell^2+O(H^2)$.
Using \eqref{eq:ar1-block-scales} and
$\sum_{\ell\le M}\ell^4=M^5/5+O(M^4)$ gives
\[
 \sum a_\ell^2=\frac{\eta^2}{5\sqrt\xi}H^{5/2}+O(H^2),\qquad
 \sum a_\ell n_\ell^2=\frac{\eta\sqrt\xi}{5}H^{5/2}+O(H^2).
\]
This proves \eqref{eq:ar1-calendar-upper}.
Its coefficient is minimized at $\xi_*=2\beta_k\eta/(\nu r)$, giving the
right side of \eqref{eq:ar1-converse-final}. Combining the two proves
\eqref{eq:ar1-main} under the stated common contract.
\end{proof}

\paragraph{Scope of the result.}
At $\rho=0$ the coefficient reduces to
$\sqrt2\sqrt{kr}\eta^{3/2}/(5\sigma^2)$, with the two-sided difference rate
$r$ rather than a one-sided healthy rate. The oracle and achieved experiment
share physical time, raw observation channel, stationary initialization and
the complete nuisance class. The result concerns a known planned horizon,
not an anytime sharp coefficient, a sequential ARL statement, a fixed finite
amplitude box, or a physically free choice of noise correlation.
The physical map from a controlled plant to this experiment and its feasible
task class requires an additional, separately stated model.
